% !Mode:: "TeX:UTF-8"
\chapter{总结与展望}

\section{本文工作总结}

本文围绕"基于EDA工具反馈的LLM Verilog生成与自动优化"这一研究问题，完成了以下工作。

在系统实现方面，本文构建了一个AutoChip风格的RTL自动生成与反馈修复原型系统，实现了完整的"生成—编译—仿真—反馈—修复"闭环，包括任务加载（支持VerilogEval和RTLLM双benchmark适配）、大语言模型统一API调用、Verilog代码提取与编译仿真执行、多粒度反馈信息构建、多种提示词策略支持，以及迭代修复循环控制。

在benchmark接入方面，本文在中期阶段使用VerilogEval-Human 20题子集完成了基础验证实验，随后引入RTLLM 2.0作为更具挑战性的评估基准。通过对RTLLM全部50个设计的Icarus Verilog兼容性审计（41/50完全可用），从中精选12道高难度、多类别覆盖的设计构成RTLLM\_STUDY\_12正式研究子集。

在实验研究方面，本文设计并执行了七组递进式实验，覆盖了基线验证、正式实验、消融实验、稳定性重复实验、反馈粒度实验、多轮对话反馈实验和提示词策略实验。在工程质量方面，建立了完善的实验产物管理和审计机制。

\section{主要实验结论}

% 审阅报告高优先级修改：从机械"结论X"列表改为自然段
基于七组实验的结果，本文得出以下主要结论。

EDA工具反馈能够提升RTL代码生成正确率。在VerilogEval-Human上，Haiku模型的通过率从80\%提升至90\%；在RTLLM\_STUDY\_12上，三个模型均观察到反馈带来的通过率提升。在更难benchmark与强模型条件下，反馈的价值更为突出：GPT-5.4从50\%提升至83\%，绝对提升33个百分点。

反馈的收益并非简单的多次重试。消融实验表明，GPT-5.4上反馈相对于retry-only的增益平均为+16.7个百分点，且存在仅反馈可解而15次独立重试均无法通过的设计题目。稳定性实验确认反馈信息约贡献了总提升的57\%，多采样贡献约43\%。

反馈效果具有模型依赖性。GPT-5.4上反馈始终优于retry-only（4/4轮均成立），而Sonnet 4.6上反馈在当前设置下反而不如retry-only（0/4轮优于retry-only）。这一结论在重复实验中完全稳定。

反馈信息量存在最佳区间。在本文实验设置下，反馈粒度实验呈现倒U型趋势：仅编译反馈（L2）和简洁反馈（L3）达到最优效果，而丰富反馈（L4）反而出现下降。多轮对话反馈和提示词策略可作为补充优化方向，其中多轮对话在控制候选数量后带来+17个百分点的提升，提示词策略方面反馈循环将所有策略拉平至相同水平（92\%）。

\section{系统局限性}

本文的系统实现和实验研究存在以下局限性。RTLLM\_STUDY\_12包含12道精选题目，虽然覆盖四大类别和较高难度梯度，但仍不能代表所有类型的RTL设计任务。本文的正式实验集中在三个模型上，其他模型是否表现出类似的反馈响应模式尚未验证。当前系统的反馈信息主要基于仿真文本输出，缺少更精确的波形级定位。部分实验为局部探索，统计效力有限。

\section{后续工作展望}

\subsection{扩大benchmark与任务规模}

可将实验从RTLLM\_STUDY\_12扩展到更大的任务子集或全量RTLLM，引入其他高难度RTL生成benchmark，进行更细粒度的分析。

\subsection{更多模型与自适应反馈策略}

纳入更多闭源和开源大语言模型进行对照实验，探索根据模型特性和错误类型动态选择反馈策略的自适应机制。

\subsection{更精细的仿真反馈与波形分析}

引入信号级差异分析，将仿真反馈从文本摘要提升到波形关键时间点定位的精度。

\subsection{检索增强与领域知识补充}

针对天花板题暴露的领域知识缺口问题，探索通过检索相似的Verilog设计实现来补充模型的先验知识\cite{hdldebugger}。

\subsection{形式化验证集成}

探索将形式化验证工具集成到反馈循环中，利用形式化工具提供的反例作为更精确的修复线索。

\section{本章小结}

本章总结了本文完成的系统实现和实验分析工作，概括了主要实验结论，诚实讨论了系统和实验的局限性，并从benchmark扩展、模型覆盖、反馈精细化、知识增强和形式化验证等方向展望了后续研究。
